\documentclass[10pt,a4paper]{amsart}
\usepackage[margin=19mm]{geometry}
\usepackage{amsmath,amssymb,mathtools}
\usepackage[scr=rsfs,cal=euler]{mathalfa}
\usepackage{xfrac}
\usepackage[unicode,hidelinks,bookmarks=false]{hyperref}
\newcommand{\ie}{i.e.\ }
\newcommand{\cf}{cf.\ }
\newcommand{\resp}{respectively\ }
\newcommand{\Subsection}{\subsection}
\providecommand{\nobreakdash}{\nobreak\hbox{-}}
\setlength{\emergencystretch}{2em}
\multlinegap0pt
\usepackage{bm}
\usepackage{bbm}
\usepackage{dsfont}
\usepackage{xspace}
\usepackage{cancel}
\usepackage{mathtools}
\usepackage{amsthm}
\makeatletter
\@ifundefined{theorem}{\newtheorem{theorem}{Theorem}}{}
\makeatother
\title[Limiter Spaces: A Universal Extension for Limits of Real Seq]{Limiter Spaces: A Universal Extension for Limits of Real Sequences}
\author{Steven Lapp, Marina Tvalavadze}
\date{Source: arXiv:2604.24569v1 (CC BY 4.0)}
\begin{document}
\maketitle

\noindent\emph{Source and licence.} Limiter Spaces: A Universal Extension for Limits of Real Sequences. Version arXiv:2604.24569v1 (CC BY 4.0), distributed under the Creative Commons Attribution 4.0 International licence, as recorded in the arXiv licence metadata for this exact version and shown on the abstract page https://arxiv.org/abs/2604.24569v1. Full rights notice: licenses/arxiv-2604.24569-CC-BY-4.0.txt.\par\smallskip

\noindent\textit{Editorial scope.} Counted unit: the Theorem whose source label is \texttt{None} in the article (source0.tex lines 440--444, source label \texttt{None}), delivered together with the article's own complete original proof of it (source0.tex lines 445--560, one \texttt{proof} environment of 116 lines). No mathematics is written, repaired or re-derived: statement and proof are the article's own text line for line. Required context retained uncounted: none. Every definition, display and label of those lines is retained unchanged. Omitted from this package and not redistributed: 1--439, 561--664. Nothing inside a retained excerpt is omitted or altered: every retained body is delivered exactly as the article's lines are written, so no omission receipt of any kind accompanies this package. Citations: the retained excerpts carry no citation command at all, so no bibliography entry of the article is transcribed and no reference is redistributed. Reference adaptation: none was needed and none was made; every reference inside the delivered unit resolves inside it, so no unresolved reference or citation is produced. Printed slips kept literal: no printed word, symbol, hypothesis or number of the counted statement or of its proof was altered. Layout and class: the article's own class is \texttt{amsart} with packages including amsmath, amssymb, babel, tikz, graphicx, anysize, amsfonts, mathrsfs, xy, overpic, fancyhdr, hyperref, ytableau, multicol; the wrapper is the lane's pinned \texttt{amsart} editorial wrapper at 10pt on A4, which restates the theorem-like environments the retained text needs and adds the article's own macro definitions that the retained text needs (none), with extra wrapper packages (bm, bbm, dsfont, xspace, cancel, mathtools). The delivered numbering is the unit's own. Assets: no retained span contains a figure environment, a graphics-inclusion command, a PDF-page inclusion, a table or a TikZ picture; no image file is copied into this package, the package declares no asset dependency and no image file is redistributed with it. Licence: version arXiv:2604.24569v1 of this article is distributed under the Creative Commons Attribution 4.0 International licence, as recorded in the arXiv licence metadata for this exact version and shown on the abstract page https://arxiv.org/abs/2604.24569v1. A full rights notice is delivered at licenses/arxiv-2604.24569-CC-BY-4.0.txt. Characters: every retained excerpt is pure ASCII, so the delivered engine, XeLaTeX with its default Latin Modern text font, reports no missing character.
\begin{theorem} For any sequence $(x_n) \in Seq(\mathbb{R})$, the sequence $(i(x_n))$  converges to $[\hspace{0.5mm}(x_n)\hspace{0.5mm}]$ in $Seq(\mathbb{R})/\hspace{-1mm}\sim.$
Furthermore, the triple $(Seq(\mathbb{R})/\hspace{-1mm}\sim,\,i,\,\varphi_{\sim})$ is the 
Limiter.  
  
\end{theorem}

\begin{proof}
 First, we show that the sequence $(i(x_n))$ converges to 
 $[\hspace{0.5mm}(x_n)\hspace{0.5mm}]$ for any sequence $(x_n) \in Seq(\mathbb{R}).$
For the sake of contradiction, let us assume that we can find a sequence $(x_n) \in Seq(\mathbb{R})$ so that $(i(x_n))$ does not converge to $[\hspace{0.5mm}(x_n)\hspace{0.5mm}]$. Then there exists a basic set $O^*$ such that $[\hspace{0.5mm}(x_n)\hspace{0.5mm}] \in O^*$
and a subsequence $(x_{n_k})$ with the property that $i(x_{n_k}) \not \in O^*$ for all $n_k \in \mathbb{N}$. 
Since a cluster set of any sequence in $\overline{\mathbb{R}}$ is non-empty, we can find a subsequence of $(x_{n_k})$ which is convergent in $\overline{\mathbb{R}}.$
Denote this subsequence by $(x_{n_{k_w}})$ and its limit by $L \in \overline{\mathbb{R}}$.
Due to being a subsequence of $(x_n)$, $L \in \Lambda(x_n) \subseteq O$ by definition of $O^*$.

Since the inclusion map $i: \overline{\mathbb{R}} \to Seq(\mathbb{R})/ \hspace{-1mm} \sim \hspace{1mm}$ is continuous and, in particular, sequentially continuous,  we have that $i(x_{n_{k_w}})$ converges to 
$i(L)$. Consequently, $i(L)$ being a member of an open set $O^*$ allow us to find a natural number $N \in \mathbb{N}$ with the property that $i(x_{n_{k_w}}) \in O^*$ whenever $n_{k_w} > N$. 
This is a contradiction. Therefore, $i(x_n)$ converges to $[ \hspace{0.5mm}(x_n) \hspace{0.5mm} ] = \varphi_{\sim}(x_n)$. This establishes the
first axiom of the Limiter. \\ 

We next verify that $(Seq(\mathbb{R})/\hspace{-1mm}\sim,\,i, \varphi_{\sim})$ satisfies the remaining axioms, that is,
\par\medskip
\hspace{5mm} (2) If $(x_n)$ is convergent in $\overline{\mathbb{R}}$, then $ \displaystyle i( \hspace{1mm} \lim_{n\to\infty} x_n \hspace{1mm}) = \varphi_{\sim}(\displaystyle \lim_{n \to \infty}(x_n)) = [\hspace{0.5mm}(\displaystyle \lim_{n \to \infty}(x_n))\hspace{0.5mm}];$

\vspace{5mm}

\hspace{5mm} (3)  $[\hspace{0.5mm}(x_n)\hspace{0.5mm}] = \varphi_{\sim}(x_n) = \varphi_{\sim}(y_n) = [\hspace{0.5mm}(y_n)\hspace{0.5mm}]$ if and only if $(x_n) \sim (y_n);$

\vspace{5mm}

\hspace{5mm} (4) For all $\alpha \in Seq(\mathbb{R})/\hspace{-1mm}\sim$, there exists $(z_n) \in Seq(\mathbb{R})$ such that $[\hspace{0.5mm} (z_n) \hspace{0.5mm}] = \varphi_{\sim}(z_n) = \alpha$;

\vspace{5mm}

\hspace{5mm} (5) For any open set $\widetilde{O}$ with $\phi_{\sim}(x_n) \in \widetilde{O}$ there exists a sub-neighborhood $W \subseteq \widetilde{O}$ of 

\vspace{3mm}

\hspace{11mm} $\varphi_{\sim}(x_n)$ such that $[(y_n)] \in W$ if and only if $\Lambda(x_n) \subseteq i^{-1}(W).$

\vspace{7mm}

In the above axioms, we denote $\underset{n\to \infty}{Lim}$ by $\varphi_{\sim}$ for convenience. We note that conditions (2)-(4) follow from the construction of $Seq(\mathbb{R})/\sim$. It remains to prove the fifth. 

Let $\widetilde{O}$ be an open set containing 
$ \varphi_{\sim}(x_n) = [\hspace{0.5mm}(x_n)\hspace{0.5mm}] \in \widetilde{O}$. Then $\widetilde{O}$ is the union of basic sets by definition of our topology. As such, one of the basic sets $W$ in the union must contain our specified equivalence class $[\hspace{0.5mm}(x_n)\hspace{0.5mm}]$. For $W$ there exists $V \subseteq [-\infty,\infty]$ such that  $W=V^{*}$. We 
have therefore found an open set $V \subseteq [-\infty,\infty]$ such that $[(x_n)] \in V^* \subseteq \widetilde{O} $. 
We recall $V^* = \{[(x_n)] \hspace{1mm} | \hspace{1mm} \Lambda(x_n) \subseteq V \}$. By our proof of $i$ being 
an embedding, $i^{-1}(V^*) = V$. Consequently, $\varphi_{\sim}(y_n) = [(y_n)] \in V^* $ if and only if
$\Lambda(y_n) \subseteq V = i^{-1}(V^*)$. This proves the final axiom. Therefore, $(Seq(\mathbb{R})/\hspace{-1mm}\sim,\,i, \varphi_{\sim})$ satisfies the axioms of being a Limiter. \\


It remains to show that it is universal. 
To prove universality, we consider an arbitrary triple $(\mathcal {Y},\,j,\, {L})$ where $j:\overline{\mathbb{R}} \to \mathcal Y$ and ${L}: Seq(\mathbb{R}) \to \mathcal Y$ satisfying axioms of the Limiter. Let $\mathcal X$ denote $Seq(\mathbb{R})/\hspace{-1mm}\sim$. Then we define $T: \mathcal X \to \mathcal Y$ by
$$T([\hspace{0.5mm}(x_n)\hspace{0.5mm}]) = {L}(x_n).$$
By axiom 3,  $T$ is not only a well defined function but also injective. Furthermore, axiom 4 implies that $T$ is surjective. It is therefore a bijection.

We will show that $T$ is the unique continuous function satisfying the following: 



$$T \circ i = j \hspace{2mm} \text{and} \hspace{2mm} T \circ \varphi_{\sim} = {L}$$ 

We begin by proving continuity of $T$.
Let $O$ be an open set in $\mathcal Y$. Then, for every point $\alpha \in O$, there exists an open subset $V_{\alpha} \subseteq O$ such that $\alpha \in V_{\alpha}$ and  for all sequences 
$(y_n) \in Seq(\mathbb{R}):$  ${L}(y_n) \in V$ if and only if  $j(\Lambda(y_n)) \subseteq V_{\alpha}$. This is a consequence of axioms 4 and 5. 
Under the same constraints, continuity of the embedding $j: \overline{\mathbb{R}} \to \mathcal{Y}$, asserts that $j^{-1}(V_{\alpha})$ is an open subset of $\overline{\mathbb{R}}$. Due to
this property, we have that $(j^{-1}(V_{\alpha}))^*$ is basic set in $\mathcal X= Seq(\mathbb{R})/\hspace{-1mm}\sim$ for any $\alpha \in O$. Taking the image of these basic sets under $T$ gives us the following: 
\[
(j^{-1}(V))^* = \{ \hspace{1mm} [(x_n)]\in \mathcal{X} \hspace{1mm} | \hspace{1mm} [(x_n)] \in  (j^{-1}(V)))^* \hspace{1mm} \} = 
\{[(x_n)]\in \mathcal{X} \hspace{1mm} | \hspace{1mm} \Lambda(x_n) \subseteq j^{-1}(V)\},
\]
\[
\begin{aligned}
T\bigl((j^{-1}(V_{\alpha}))^*\bigr) 
  &= \{\, T([x_n]) \;\big|\; \Lambda(x_n) \subseteq j^{-1}(V_{\alpha}) \,\} \\[4pt]
  &= \{\, {L}(x_n) \;\big|\; j(\Lambda(x_n)) \subseteq V_{\alpha} \,\} \\[4pt]
  &= \{\, {L}(x_n) \;\big|\; {L}(x_n) \in V_{\alpha }\,\}= V_{\alpha}
\end{aligned}
\]



This shows  that $T\bigl((j^{-1}(V_{\alpha}))^*\bigr) = V_{\alpha}$ for any $\alpha \in O$. Due to $T$ being a bijection, the above is equivalent to $T^{-1}(V_{\alpha}) = (j^{-1}(V_{\alpha}))^{*}$ for any $\alpha \in O$. It follows that 


$$\displaystyle T^{-1}(O) = T^{-1} (\bigcup_{\alpha \in O} V_{\alpha}) = \bigcup_{\alpha \in O} T^{-1}(V_{\alpha}) = \bigcup_{\alpha \in O} (j^{-1}(V_{\alpha}))^*$$
 


This proves that $T^{-1}(O)$ can be represented as the union of basic sets, proving continuity of $T$.


We now prove that $T \circ i = j$ and $T \circ \phi_{\sim} = {L}.$
Consider $x \in \overline{\mathbb{R}}$. Then, by construction of $i$ and 
Theorem 1, we have that $\underset{n\to \infty}{Lim}(x_n) = x$ for all sequence in the equivalence class $i(x)$.
Consequently, the following are equal by axiom 2:



$$T(i(x))= T([(x)])={L}((x))=j(\hspace{0.5mm} \underset{n\to \infty}{lim}(x) \hspace{0.5mm}) = j(x)$$


Due to the above argument holding for all $x \in \overline{\mathbb{R}}$, we conclude that $T \circ i  = j$. Similarly, 
we can conclude that $T \circ \varphi_{\sim} = {L}$

$$T( \hspace{0.5mm}\phi_{\sim}(x_n)\hspace{0.5mm}) = T(\hspace{0.5mm}[\hspace{0.5mm}(x_n)\hspace{0.5mm}]\hspace{0.5mm}) = {L}(x_n)$$


\noindent Therefore,  $T \circ i = j$ and $T \circ \varphi_{\sim} = {L}$. 

To conclude the proof of universality, we show 
that $T$ is unique.

Let $T': \mathcal{X}  \to \mathcal{Y}$ be a continuous function satisfying our conditions. Then, for
any equivalence class $[(x_n)] \in \mathcal X$, the following holds
$$T'(\hspace{0.5mm}[\hspace{0.5mm}(x_n)\hspace{0.5mm}]\hspace{0.5mm}) = T'(\hspace{0.5mm} \varphi_{\sim}(x_n) \hspace{0.5mm} ) = (T'\circ \varphi_{\sim})(x_n) =  {L}(x_n) = T(\hspace{0.5mm}[\hspace{0.5mm}(x_n)\hspace{0.5mm}]\hspace{0.5mm}).$$
This proves 
uniqueness of the function $T$ and in turn universality. 

We have successfully constructed the Limiter of the real numbers. 
\end{proof}

\end{document}
